\chapter*{Abstract}
\addcontentsline{toc}{chapter}{Abstract}

Programming assessments check who a student is at the start of a session,
and whether the submitted code was copied or generated at the end. Nothing
checks that the same person kept writing in between. This thesis studies
\emph{authorship continuity}: detecting when a different person, such as a
friend or a paid proxy, takes over partway through a coding session. We
build Scribelab, a lab code editor that records keystroke timing, pastes,
edits and code runs without any webcam or screen recording, and use it in
real lab classes and competitions at IIIT Surat. Each student's earliest ten
active minutes form a personal baseline; every later 30-second window is
scored against it in a rhythm, a key-pair and a way-of-working channel, and a
CUSUM change-point detector flags a likely takeover. Teachers see an evidence
timeline, and a flag leads only to a conversation, never a penalty.

The evaluation uses classroom data only (1,082 recorded sessions of 276 students) and
builds labelled takeovers by splicing two students' real sessions on the same
question. With thresholds frozen on development students' real sessions for
0.2 false alarms per active hour, the detector catches 31\% of spliced
takeovers (95\% interval 18--45\%) with a median delay of ten minutes, and the
false-alarm rate held on unseen students (0.15 per hour); a plain per-key
timing baseline does about as well. Thresholds calibrated on simulated typists
raise eight times the false alarms on real classwork. Typing rhythm is a stable
personal trait over three weeks, while the way of working drifts with the task
and agrees only weakly with signs of copying (AUC 0.61--0.74 for outside
pastes); rhythm and Nosi's transcription measure do not. If one session in a
hundred held a takeover, about one flag in nineteen would be real. A flag can
start a conversation, and nothing more.

\vspace{1em}
\noindent\textbf{Keywords:} keystroke dynamics, authorship continuity,
change-point detection, CUSUM, programming education, academic integrity.
